\begin{table}[!h]
\centering
\caption{\label{tab:table_balance}Descriptive statistics -- treatment vs. control municipalities, national sample (Cohort B, 2014 baseline)}
\centering
\fontsize{9}{11}\selectfont
\begin{threeparttable}
\begin{tabular}[t]{lllll}
\toprule
Variable & Treatment
mean (SD) & Control
mean (SD) & Diff. (T\$-\$C) & p\\
\midrule
\textbf{Municipalities (N)} & \textbf{637} & \textbf{8284} & \textbf{} & \textbf{}\\
Registered voters (2014) & 420.14 (544.38) & 503.11 (609.48) & -82.97*** & 0.000\\
Median household income (EUR, 2014) & 19280.11 (1999.91) & 19610.86 (2124.99) & -330.75*** & 0.000\\
Municipal investment per capita (EUR, 2013) & 9892.13 (5594.35) & 9644.71 (5219.01) & +247.42 & 0.280\\
Wind speed at 100m (m/s) & 6.88 (0.61) & 6.82 (0.68) & +0.07** & 0.011\\
Incumbent candidacy rate (\%, 2014) & 70.99 (45.42) & 72.02 (44.89) & -1.03 & 0.597\\
Incumbent reelection rate (\%, 2014) & 61.85 (48.61) & 61.75 (48.60) & +0.10 & 0.960\\
Winner vote share (\% valid votes, 2014) & 81.61 (14.95) & 81.48 (15.75) & +0.14 & 0.826\\
Abstention rate (\% registered, 2014) & 23.16 (8.72) & 24.52 (8.63) & -1.36*** & 0.000\\
\bottomrule
\end{tabular}
\begin{tablenotes}
\item \textit{Note: } 
\item TREATMENT: rural density 5-7, 0 wind farms in 2014 -> >=1 by 2020 (Cohort B). CONTROL: rural density 5-7, never treated, department with > 3 treated municipalities, probit-matched on wind speed + log(pop) + log(income) + log(investment) + department FE (P10 threshold of treated). All variables measured at 2014 (pre-treatment baseline). Diff. = Treatment mean - Control mean, two-sample t-test. *** p<0.01, ** p<0.05, * p<0.10.
\end{tablenotes}
\end{threeparttable}
\end{table}
